% GENERATED FILE. Edit canonical lessons, then run npm run generate:books.
% canonical-source-hash: fnv1a64:c97b39afd5224779
% canonical-lessons: HI-C350-udaharan, HI-C350-maslan, HI-C350-yani, HI-C350-spast, HI-C350-ginana, HI-R350-review-343-350

\chapter{For Example}
\label{ch:hi-for-example}
\hlchaptermodality{350}

\begin{chapteropening}
\textbf{By the end of this chapter:} \emph{I can back up a point with an example.}

Complete the last lesson of chapter 350: I can back up a point with an example.
\end{chapteropening}

\section[\texorpdfstring{\hi{उदाहरण} (udāharaṇ)}{उदाहरण (udāharaṇ)}]{\hi{उदाहरण} (udāharaṇ) — an example}

\label{lesson:HI-C350-udaharan}



\begin{quote}
Before the new one: say the Hindi for to admit, then the Hindi for a misunderstanding.
\end{quote}

\subsection*{You'll want to know: \hi{उदाहरण}}
\textbf{\hi{उदाहरण}} — \emph{udāharaṇ} — \textquotedblleft{}an example\textquotedblright{}.

\begin{quote}
\hi{शिक्षक} \hi{ने} \hi{एक} \hi{उदाहरण} \hi{दिया।}
\end{quote}

\begin{quote}
\emph{śikṣak ne ek udāharaṇ diyā.}
\end{quote}

\begin{quote}
The teacher gave an example.
\end{quote}

Masculine.

\begin{cousinweb}
\textbf{\hi{उदाहरण}} (\emph{udāharaṇ}) is a tatsama, borrowed from Sanskrit \emph{udāharaṇa}, \textquotedblleft{}an example, an illustration\textquotedblright{} (Monier-Williams).
\end{cousinweb}

\begin{grammarlens}[title={What you've built}]
One more word for giving reasons and opinions.
\end{grammarlens}

\subsection*{Your turn}
\begin{itemize}
\raggedright
  \item \emph{Say it:} \emph{udāharaṇ}
  \item \emph{Say it:} \emph{udāharaṇ}, once more
  \item \emph{Say it:} \emph{śikṣak ne ek udāharaṇ diyā.} — The teacher gave an example.
  \item \emph{Recall:} say the Hindi for to admit, then the Hindi for a misunderstanding, then say \emph{udāharaṇ} again
\end{itemize}

\begin{tcolorbox}[breakable,colback=teal!4,colframe=teal!35!black,title={Before you move on}]
What does \hi{उदाहरण} mean? (\textquotedblleft{}An example\textquotedblright{}.) Say it once more.
\end{tcolorbox}
\section[\texorpdfstring{\hi{मसलन} (maslan)}{मसलन (maslan)}]{\hi{मसलन} (maslan) — for instance}

\label{lesson:HI-C350-maslan}



\begin{quote}
Before the new one: say the Hindi for a misunderstanding, then the Hindi for an example.
\end{quote}

\subsection*{You'll want to know: \hi{मसलन}}
\textbf{\hi{मसलन}} — \emph{maslan} — \textquotedblleft{}for instance\textquotedblright{}.

\begin{quote}
\hi{मुझे} \hi{फल} \hi{पसंद} \hi{हैं}, \hi{मसलन} \hi{आम} \hi{और} \hi{केला।}
\end{quote}

\begin{quote}
\emph{mujhe phal pasand haiṁ, maslan ām aur kelā.}
\end{quote}

\begin{quote}
I like fruit, for instance mangoes and bananas.
\end{quote}

\begin{cousinweb}
\textbf{\hi{मसलन}} (\emph{maslan}) is usually traced to Arabic \emph{mathalan}, \textquotedblleft{}for example\textquotedblright{}, which reached Hindi through Persian.
\end{cousinweb}

\begin{grammarlens}[title={What you've built}]
One more word for giving reasons and opinions.
\end{grammarlens}

\subsection*{Your turn}
\begin{itemize}
\raggedright
  \item \emph{Say it:} \emph{maslan}
  \item \emph{Say it:} \emph{maslan}, once more
  \item \emph{Say it:} \emph{mujhe phal pasand haiṁ, maslan ām aur kelā.} — I like fruit, for instance mangoes and bananas.
  \item \emph{Recall:} say the Hindi for a misunderstanding, then the Hindi for an example, then say \emph{maslan} again
\end{itemize}

\begin{tcolorbox}[breakable,colback=teal!4,colframe=teal!35!black,title={Before you move on}]
What does \hi{मसलन} mean? (\textquotedblleft{}For instance\textquotedblright{}.) Say it once more.
\end{tcolorbox}
\section[\texorpdfstring{\hi{यानी} (yānī)}{यानी (yānī)}]{\hi{यानी} (yānī) — that is to say, in other words}

\label{lesson:HI-C350-yani}



\begin{quote}
Before the new one: say the Hindi for an example, then the Hindi for for instance.
\end{quote}

\subsection*{You'll want to know: \hi{यानी}}
\textbf{\hi{यानी}} — \emph{yānī} — \textquotedblleft{}that is to say, in other words\textquotedblright{}.

\begin{quote}
\hi{वह} \hi{परसों} \hi{आएगा}, \hi{यानी} \hi{बुधवार} \hi{को।}
\end{quote}

\begin{quote}
\emph{vah parsoṁ āegā, yānī budhvār ko.}
\end{quote}

\begin{quote}
He'll come the day after tomorrow, that is, on Wednesday.
\end{quote}

\begin{cousinweb}
\textbf{\hi{यानी}} (\emph{yānī}) is usually traced to Arabic \emph{yaʿnī}, \textquotedblleft{}it means\textquotedblright{}, which reached Hindi through Persian.
\end{cousinweb}

\begin{grammarlens}[title={What you've built}]
One more word for giving reasons and opinions.
\end{grammarlens}

\subsection*{Your turn}
\begin{itemize}
\raggedright
  \item \emph{Say it:} \emph{yānī}
  \item \emph{Say it:} \emph{yānī}, once more
  \item \emph{Say it:} \emph{vah parsoṁ āegā, yānī budhvār ko.} — He'll come the day after tomorrow, that is, on Wednesday.
  \item \emph{Recall:} say the Hindi for an example, then the Hindi for for instance, then say \emph{yānī} again
\end{itemize}

\begin{tcolorbox}[breakable,colback=teal!4,colframe=teal!35!black,title={Before you move on}]
What does \hi{यानी} mean? (\textquotedblleft{}That is to say, in other words\textquotedblright{}.) Say it once more.
\end{tcolorbox}
\section[\texorpdfstring{\hi{स्पष्ट} (spaṣṭ)}{स्पष्ट (spaṣṭ)}]{\hi{स्पष्ट} (spaṣṭ) — clear, explicit}

\label{lesson:HI-C350-spast}



\begin{quote}
Before the new one: say the Hindi for for instance, then the Hindi for that is to say.
\end{quote}

\subsection*{You'll want to know: \hi{स्पष्ट}}
\textbf{\hi{स्पष्ट}} — \emph{spaṣṭ} — \textquotedblleft{}clear, explicit\textquotedblright{}.

\begin{quote}
\hi{नियम} \hi{बिलकुल} \hi{स्पष्ट} \hi{है।}
\end{quote}

\begin{quote}
\emph{niyam bilkul spaṣṭ hai.}
\end{quote}

\begin{quote}
This rule is perfectly clear.
\end{quote}

\begin{cousinweb}
\textbf{\hi{स्पष्ट}} (\emph{spaṣṭ}) is a tatsama, borrowed from Sanskrit \emph{spaṣṭa}, \textquotedblleft{}plainly seen, evident\textquotedblright{} (Monier-Williams).
\end{cousinweb}

\begin{grammarlens}[title={What you've built}]
One more word for giving reasons and opinions.
\end{grammarlens}

\subsection*{Your turn}
\begin{itemize}
\raggedright
  \item \emph{Say it:} \emph{spaṣṭ}
  \item \emph{Say it:} \emph{spaṣṭ}, once more
  \item \emph{Say it:} \emph{niyam bilkul spaṣṭ hai.} — This rule is perfectly clear.
  \item \emph{Recall:} say the Hindi for for instance, then the Hindi for that is to say, then say \emph{spaṣṭ} again
\end{itemize}

\begin{tcolorbox}[breakable,colback=teal!4,colframe=teal!35!black,title={Before you move on}]
What does \hi{स्पष्ट} mean? (\textquotedblleft{}Clear, explicit\textquotedblright{}.) Say it once more.
\end{tcolorbox}
\section[\texorpdfstring{\hi{गिनाना} (ginānā)}{गिनाना (ginānā)}]{\hi{गिनाना} (ginānā) — to list, to enumerate (reasons, faults)}

\label{lesson:HI-C350-ginana}



\begin{quote}
Before the new one: say the Hindi for that is to say, then the Hindi for explicit.
\end{quote}

\subsection*{You'll want to know: \hi{गिनाना}}
\textbf{\hi{गिनाना}} — \emph{ginānā} — \textquotedblleft{}to list, to enumerate (reasons, faults)\textquotedblright{}.

\begin{quote}
\hi{मीरा} \hi{ने} \hi{तीन} \hi{कारण} \hi{गिनाए।}
\end{quote}

\begin{quote}
\emph{Mīrā ne tīn kāraṇ gināe.}
\end{quote}

\begin{quote}
Mira listed three reasons.
\end{quote}

Past, with \hi{ने} (he or she …): \hi{उसने} \hi{कई} \hi{कारण} \hi{गिनाए}, \emph{usne kaī kāraṇ gināe}.

\begin{cousinweb}
\textbf{\hi{गिनाना}} (\emph{ginānā}) is the causative of \hi{गिनना}, \textquotedblleft{}to count\textquotedblright{}, which is inherited from Sanskrit \emph{gaṇayati}, \textquotedblleft{}counts\textquotedblright{} (CDIAL 3993): to make the items count out, one by one.
\end{cousinweb}

\begin{grammarlens}[title={What you've built}]
One more word for giving reasons and opinions.
\end{grammarlens}

\subsection*{Your turn}
\begin{itemize}
\raggedright
  \item \emph{Say it:} \emph{ginānā}
  \item \emph{Say it:} \emph{ginānā}, once more
  \item \emph{Say it:} \emph{Mīrā ne tīn kāraṇ gināe.} — Mira listed three reasons.
  \item \emph{Recall:} say the Hindi for that is to say, then the Hindi for explicit, then say \emph{ginānā} again
\end{itemize}

\begin{tcolorbox}[breakable,colback=teal!4,colframe=teal!35!black,title={Before you move on}]
What does \hi{गिनाना} mean? (\textquotedblleft{}To list, to enumerate (reasons, faults)\textquotedblright{}.) Say it once more.
\end{tcolorbox}
\section[(dialogue)]{Looking back: from What Caused It to For Example}

\label{lesson:HI-R350-review-343-350}



\begin{quote}
Say the Hindi for since and an effect.
\end{quote}

\subsection*{Your turn}
Say these aloud:

\begin{itemize}
\raggedright
  \item since, an effect, to be joined, a basis, grounds, responsible
  \item a consequence, an influence, to rise, to make worse, a source
  \item a purpose, a motive, for the sake of, a means, an intention
  \item an argument, a line of reasoning, to quarrel, to interrupt, but rather
  \item a side, against, to oppose, to support, to agree
  \item proof (formal), to prove, a fact, conclusive, an accusation
  \item to blurt out, to deny, to confirm, to admit, a misunderstanding
  \item an example, for instance, that is to say, explicit, to list
\end{itemize}

\begin{tcolorbox}[breakable,colback=teal!4,colframe=teal!35!black,title={Before you move on}]
Since? (\textbf{\hi{चूँकि}}, \emph{cū̃ki}.) An effect? (\textbf{\hi{असर}}, \emph{asar}.) To be joined? (\textbf{\hi{जुड़ना}}, \emph{juṛnā}.) A basis, grounds? (\textbf{\hi{आधार}}, \emph{ādhār}.) Responsible? (\textbf{\hi{ज़िम्मेदार}}, \emph{zimmedār}.) A consequence? (\textbf{\hi{परिणाम}}, \emph{pariṇām}.) An influence? (\textbf{\hi{प्रभाव}}, \emph{prabhāv}.) To rise? (\textbf{\hi{बढ़ना}}, \emph{baṛhnā}.) To make worse? (\textbf{\hi{बिगाड़ना}}, \emph{bigāṛnā}.) A source? (\textbf{\hi{स्रोत}}, \emph{srot}.) A purpose? (\textbf{\hi{उद्देश्य}}, \emph{uddeśya}.) A motive? (\textbf{\hi{मक़सद}}, \emph{maqsad}.) For the sake of? (\textbf{\hi{ख़ातिर}}, \emph{khātir}.) A means? (\textbf{\hi{साधन}}, \emph{sādhan}.) An intention? (\textbf{\hi{इरादा}}, \emph{irādā}.) An argument? (\textbf{\hi{बहस}}, \emph{bahas}.) A line of reasoning? (\textbf{\hi{तर्क}}, \emph{tark}.) To quarrel? (\textbf{\hi{झगड़ना}}, \emph{jhagaṛnā}.) To interrupt? (\textbf{\hi{टोकना}}, \emph{ṭoknā}.) But rather? (\textbf{\hi{बल्कि}}, \emph{balki}.) A side? (\textbf{\hi{पक्ष}}, \emph{pakṣ}.) Against? (\textbf{\hi{ख़िलाफ़}}, \emph{khilāf}.) To oppose? (\textbf{\hi{विरोध} \hi{करना}}, \emph{virodh karnā}.) To support? (\textbf{\hi{समर्थन} \hi{करना}}, \emph{samarthan karnā}.) To agree? (\textbf{\hi{सहमत} \hi{होना}}, \emph{sahmat honā}.) Proof (formal)? (\textbf{\hi{प्रमाण}}, \emph{pramāṇ}.) To prove? (\textbf{\hi{साबित} \hi{करना}}, \emph{sābit karnā}.) A fact? (\textbf{\hi{तथ्य}}, \emph{tathya}.) Conclusive? (\textbf{\hi{पुख़्ता}}, \emph{pukhtā}.) An accusation? (\textbf{\hi{आरोप}}, \emph{ārop}.) To blurt out? (\textbf{\hi{उगलना}}, \emph{ugalnā}.) To deny? (\textbf{\hi{इनकार} \hi{करना}}, \emph{inkār karnā}.) To confirm? (\textbf{\hi{पुष्टि} \hi{करना}}, \emph{puṣṭi karnā}.) To admit? (\textbf{\hi{स्वीकार} \hi{करना}}, \emph{svīkār karnā}.) A misunderstanding? (\textbf{\hi{ग़लतफ़हमी}}, \emph{ġalatfahmī}.) An example? (\textbf{\hi{उदाहरण}}, \emph{udāharaṇ}.) For instance? (\textbf{\hi{मसलन}}, \emph{maslan}.) That is to say? (\textbf{\hi{यानी}}, \emph{yānī}.) Explicit? (\textbf{\hi{स्पष्ट}}, \emph{spaṣṭ}.) To list? (\textbf{\hi{गिनाना}}, \emph{ginānā}.)
\end{tcolorbox}
